\chapter{Deep Dive: Synchronous Calls --- OpenFeign, Resilience4j, and Cascading Failure}
\label{ch:dd-resilience}

Chapter~6 drew the line between the two ways services talk: HTTP when
the caller needs an answer now, events when it needs a fact recorded.
This chapter stays on the HTTP side and asks the question that every
distributed system eventually asks at 3\,a.m.: when the service I am
calling is slow, what happens to \emph{me}? The project has the right
libraries on its classpath --- OpenFeign for the call, Resilience4j for
the protection --- and, as you will see, has wired them in a way that
protects nothing. That makes it a better teaching case than a finished
one would be.

Three questions organize the chapter:

\begin{enumerate}
  \item \textbf{What does a synchronous call cost} beyond its latency ---
  and why is a slow dependency worse than a dead one?
  \item \textbf{How does a circuit breaker decide}, with the project's
  actual numbers, and how do timeouts, retries and breakers interact?
  \item \textbf{What should a fallback return}, and when should the call
  not exist at all?
\end{enumerate}

\section{Why Feign, and what it solves}

course-service occasionally needs to know who a user is --- a name to
print on a submission, a role to check. auth-service owns that data. The
options are: call auth-service over HTTP, copy the data into
course-service via events, or put both in one database. The project
keeps the databases separate (Chapter~7) and reserves events for facts
that happened; a lookup by id is a question, so it is a call.

OpenFeign turns that call into an interface:

\begin{lstlisting}[language=Java]
@FeignClient(name = "auth-service",
             fallback = AuthServiceFallback.class)          (*@\co{1}@*)
public interface AuthServiceClient {

    @GetMapping("/users/{id}")                              (*@\co{2}@*)
    UserDto getUserById(@PathVariable("id") Long id);        (*@\co{3}@*)
}
\end{lstlisting}

\begin{description}
  \item[\co{1}] \code{name} is the Eureka application name. Feign asks
  the discovery client for instances of \code{auth-service} and
  load-balances across them on the client side --- no gateway, no
  hard-coded host. \code{fallback} names the class to call instead when
  the call fails. Hold that thought.
  \item[\co{2}] The same Spring MVC annotations as a controller, read in
  reverse: this is what to \emph{send}. The path is the service's own
  (\code{/users/\{id\}}), not the gateway's (\code{/api/auth/users}),
  because service-to-service calls bypass the gateway.
  \item[\co{3}] The return type is a record from \code{common}. That
  record is the \emph{contract}: auth-service must produce JSON that
  deserializes into it.
\end{description}

What Feign solves is boilerplate and coupling to addresses. What it does
not solve is the one property that matters under load: every call
holds a Tomcat thread in course-service for as long as auth-service
takes to answer.

\begin{decision}[a call, not a copy --- for now]
For a by-id lookup that is rare and whose answer changes rarely, a call
is the simplest correct thing, and Feign makes it two lines. The
decision flips when any of three things becomes true: the lookup sits
on a hot path (every course listing needs a teacher's name), the data
is needed while auth-service is down (grading must not stop because
login is broken), or the shape of \code{UserDto} starts changing with
auth-service's releases. Then course-service should keep its own copy
of the fields it needs, fed by the events auth-service already
publishes (\code{ts.user.registered}, \code{ts.user.activated}) --- the
pattern Chapter~7 already applies to students.
\end{decision}

\section{The contract in \code{common}, and its versioning risk}

\code{UserDto} is a record with five fields: \code{id}, \code{email},
\code{firstName}, \code{lastName}, \code{role}. auth-service's
\code{UserController} does not return it --- it returns its own
\code{UserResponse}. The two classes agree today by coincidence of
field names, and Jackson's default of ignoring unknown properties hides
the gap: auth-service can add fields freely, and course-service will not
notice. It \emph{will} notice a rename. Renaming \code{firstName} to
\code{givenName} in \code{UserResponse} compiles on both sides, deploys
on both sides, and leaves course-service reading \code{null} names ---
no exception, no log line. The interview version: ``how do you version
an internal API?'' and the answer is: additive changes only, never
rename or remove; the consumer's DTO ignores unknowns; a consumer-driven
contract test (Spring Cloud Contract, Pact) in auth-service's build that
fails when \code{UserResponse} stops satisfying \code{UserDto}. A URL
version (\code{/v2/users}) is for breaking changes you cannot avoid,
and it means running both for as long as any consumer lags.

\section{The circuit breaker, with the project's numbers}
\label{sec:cb-math}

Resilience4j's breaker is a state machine over a \emph{sliding window}
of recent call outcomes. The project's configuration in
\code{course-service.yml}:

\begin{lstlisting}[language=yaml]
resilience4j:
  circuitbreaker:
    instances:
      authService:
        sliding-window-size: 10                    (*@\co{1}@*)
        failure-rate-threshold: 50                 (*@\co{2}@*)
        wait-duration-in-open-state: 10s           (*@\co{3}@*)
        permitted-number-of-calls-in-half-open-state: 3  (*@\co{4}@*)
\end{lstlisting}

\begin{figure}[htb]
\centering
\begin{tikzpicture}[node distance=13mm and 14mm]
  \node[boxfill] (closed) {CLOSED\\calls pass, outcomes recorded};
  \node[boxdark, right=of closed] (open) {OPEN\\calls rejected instantly};
  \node[box, below=of $(closed)!0.5!(open)$] (half) {HALF\_OPEN\\3 trial calls};
  \draw[flow] (closed) -- node[lbl,above]{$\ge$5 of 10 fail} (open);
  \draw[flow] (open.east) -- ++(6mm,0) |- node[lbl,above,pos=0.75]{after 10\,s} (half.east);
  \draw[flow] (half.north west) -- node[lbl,above,sloped]{trials ok} (closed.south);
  \draw[flowdash] (half.north east) -- node[lbl,above,sloped]{trials fail} (open.south);
\end{tikzpicture}
\caption{The breaker's three states with this project's thresholds.
Rejected calls in OPEN go to the fallback without touching the network.}
\label{fig:cb-states}
\end{figure}

\begin{description}
  \item[\co{1}] The window holds the last 10 outcomes (count-based;
  the alternative is time-based, ``the last 60 seconds''). The breaker
  does not evaluate until the window has
  \code{minimum-number-of-calls} outcomes --- unset here, so it
  defaults to the window size. Translation: the first 9 calls after
  startup can \emph{all} fail and the breaker stays CLOSED.
  \item[\co{2}] At 10 outcomes, if 5 or more failed, the breaker opens.
  A ``failure'' is an exception the breaker does not ignore. Note what
  is \emph{not} configured: \code{slow-call-duration-threshold}. Its
  default is 60\,s, so a call that takes 30\,s is a \emph{success} to
  this breaker. Slow dependencies --- the common case --- never trip
  it.
  \item[\co{3}] OPEN for 10\,s: every call is rejected with
  \code{CallNotPermittedException} in microseconds and routed to the
  fallback. auth-service gets a 10-second breather.
  \item[\co{4}] Then HALF\_OPEN: 3 calls are let through as probes. If
  their failure rate is $\ge$ 50\,\% (2 of 3), back to OPEN for another
  10\,s; otherwise CLOSED with a fresh window.
\end{description}

Now the part that makes this chapter honest. That \code{authService}
instance is bound to \emph{nothing}. No method carries
\code{@CircuitBreaker(name = "authService")}, and Feign's own
integration --- which would wrap \code{getUserById} in a breaker named
\code{AuthServiceClient\#getUserById(Long)} --- is off, because the
switch \code{openfeign.circuitbreaker.enabled} defaults to
\code{false}, and the starter that provides the
\code{CircuitBreakerFactory} it needs (the
\code{circuitbreaker-resilience4j} starter) is not on the classpath.
With that switch off, the \code{fallback} attribute in
\co{1} is silently ignored. And \code{AuthServiceClient} is not called
from any service class today. The project has a breaker configured for
a call that does not happen, protecting nothing, with a fallback that
would not run. It is a perfect specimen of how resilience is usually
misconfigured in production: every piece present, no piece connected.

\section{Timeouts stack}

When the call is wired properly, three timers surround it, and they
must be ordered or the outer ones are decorative.

\begin{itemize}
  \item \textbf{Feign's socket timeouts.} \code{connectTimeout} and
  \code{readTimeout} on the HTTP client. Unset in the project, so the
  defaults apply: 10\,s to connect, \textbf{60\,s} to read. A hung
  auth-service pins a course-service thread for a minute.
  \item \textbf{Resilience4j's TimeLimiter.} A deadline on the whole
  decorated call, independent of the socket. It only applies to
  asynchronous return types (\code{CompletableFuture}), which is why
  synchronous Feign calls rely on the socket timeout instead.
  \item \textbf{The caller's deadline.} The gateway's
  \code{response-timeout} (Chapter~31), and behind it the browser's
  patience. If the gateway gives up at 5\,s while Feign waits 60\,s, the
  user sees an error at 5\,s and course-service keeps a thread busy for
  55 more seconds serving nobody.
\end{itemize}

The rule: every hop's timeout is \emph{shorter} than the hop before it.
Browser 10\,s $>$ gateway 5\,s $>$ course-service's Feign read 2\,s. A
timeout that fires inside an already-abandoned request is wasted work;
one that fires later than its caller's is a leak.

\section{Retry: only for calls you can repeat}

Resilience4j's \code{@Retry} re-issues a failed call. On a \code{GET
/users/\{id\}} that is harmless: reading twice is the same as reading
once. On a \code{POST /enrollments} it is a duplicate enrollment,
because a timeout does not tell you whether the first request was
processed. Retry is therefore for \emph{idempotent} operations only,
or for operations made idempotent with a client-supplied key the server
deduplicates on. Two more rules from production: retry with
\emph{exponential backoff and jitter}, or every client retries at the
same instant and the second wave is bigger than the first (the ``retry
storm''); and retry \emph{inside} the breaker, so that an open breaker
stops retries from being attempted at all. Resilience4j's default aspect
order --- Retry outermost, then CircuitBreaker, then TimeLimiter, then
Bulkhead --- gets this right: a retry goes through the breaker each
time, and a rejected call is not retried.

\section{Bulkhead: bounding the blast}

A bulkhead caps how many concurrent calls may be in flight to one
dependency. Semaphore bulkhead: a counter; the 26th concurrent call is
rejected immediately. Thread-pool bulkhead: calls run on a dedicated
pool so the caller's threads are never the ones waiting --- necessary
for a servlet app whose Tomcat pool is finite. course-service runs
Tomcat with 200 threads. Without a bulkhead, 200 slow auth-service calls
consume every one of them, and course-service cannot answer a health
probe, let alone a course listing. With a semaphore bulkhead of 25, at
most 25 threads are ever waiting on auth-service and 175 remain for
everything else. That is the isolation the chapter's title promises.

\section{What the fallback returns}

\begin{lstlisting}[language=Java]
@Component
public class AuthServiceFallback implements AuthServiceClient {
    @Override
    public UserDto getUserById(Long id) {
        return null;                                       (*@\co{1}@*)
    }
}
\end{lstlisting}

\begin{description}
  \item[\co{1}] A \code{null} that the caller must remember to check.
  Every consumer of \code{getUserById} now has a hidden branch, and a
  forgotten check is a \code{NullPointerException} that fires only when
  auth-service is down --- exactly when you least want a new bug.
\end{description}

A fallback has four honest options, in order of preference: a
\emph{cached} value from the last successful call (stale but real); a
\emph{degraded} value that the UI can show (``Unknown teacher'') and
that is typed as such; an \emph{empty} result where the domain allows
one; or a \emph{typed error} the caller can map to 503 with
\code{Retry-After}. \code{null} is none of these. It is the absence of a
decision, pushed onto every caller. The fallback also needs the
\emph{cause}: a two-argument fallback method (\code{Throwable last})
can tell a breaker rejection from a 404, and a 404 --- the user does not
exist --- must \emph{not} be turned into ``try later''.

\section{Pros and cons, honestly}

\begin{center}
\small
\begin{tabular}{@{}p{0.46\linewidth}p{0.46\linewidth}@{}}
\toprule
What the project has & What it is missing \\
\midrule
Declarative client, Eureka-resolved & The client is never called \\
A breaker instance with sane thresholds & Nothing binds to it; Feign's breaker switch is off \\
A fallback class registered & Ignored; and it returns \code{null} \\
Resilience4j on the classpath & No timeouts, no retry, no bulkhead, no slow-call threshold \\
Actuator exposed & Breaker endpoints not in the exposure list \\
\bottomrule
\end{tabular}
\end{center}

And one more gap that is not about resilience at all: auth-service's
\code{SecurityConfig} permits only the public paths; \code{/users/\{id\}}
requires an authenticated context, which its
\code{GatewayHeaderAuthFilter} builds from \code{X-User-Id}. Feign sends
no such header. The call, once wired, would return 401 on every attempt
--- and a 401 is an exception, so ten of them would open the breaker
against a perfectly healthy service.

\section{What breaks}

\begin{pitfall}[the fallback that never runs]
Symptom: auth-service is stopped, and instead of the fallback's
\code{null} course-service throws \code{FeignException} with a stack
trace through the load balancer, and the request returns 500. Cause:
\code{@FeignClient(fallback = ...)} is only honored when
\code{spring.cloud.openfeign.circuitbreaker.enabled=true} \emph{and} a
\code{CircuitBreakerFactory} bean exists, which requires the
\code{spring-cloud-starter-circuitbreaker-resilience4j} dependency. With
either missing, Spring logs nothing and the attribute is inert. Fix in
Section~\ref{sec:res-enhance}.
\end{pitfall}

\begin{pitfall}[healthy service, open breaker]
Symptom: after a deploy of auth-service, course-service's breaker opens
and stays open, cycling OPEN $\to$ HALF\_OPEN $\to$ OPEN, though
auth-service answers \code{curl} fine. Cause: every call gets 401
(missing identity headers), or 404 for a deleted user, and both count
as failures. Fix: propagate identity on service-to-service calls (a
\code{RequestInterceptor}, below), and tell the breaker which exceptions
are \emph{not} dependency failures with \code{ignore-exceptions} ---
client errors mean the request was wrong, not that the service is
down.
\end{pitfall}

\begin{pitfall}[the 60-second thread]
Symptom: auth-service hangs (deadlock, stuck GC); course-service's
CPU is idle, yet it stops answering \emph{all} requests, and its
readiness probe fails so Kubernetes stops routing to it. Cause:
Feign's default 60\,s read timeout multiplied by Tomcat's 200 threads.
Each hung lookup holds a thread for a minute; with a modest request
rate every thread is waiting within seconds. The breaker cannot help
--- with the default slow-call threshold, hung calls are not failures
until they time out. Fix: a 2\,s read timeout, a slow-call threshold,
and a bulkhead of 25.
\end{pitfall}

\section{What breaks at 3\,a.m.}

auth-service's Postgres begins a long vacuum; \code{/users/\{id\}} goes
from 20\,ms to 2\,s. Assume the Feign call has been wired into the
course-listing endpoint, as the project's design invites, and nothing
else is changed.

\begin{enumerate}
  \item Each course listing now takes 2\,s instead of 20\,ms. Tomcat
  threads in course-service are busy 100$\times$ longer per request.
  \item At 50 listings per second, 100 threads are waiting on
  auth-service at any moment. At 100 per second, all 200 are, and
  requests queue in Tomcat's accept backlog.
  \item course-service's readiness probe (\code{/actuator/health} is
  not even exposed here --- Chapter~17 uses a TCP probe) keeps passing,
  because the socket accepts. The gateway keeps sending traffic.
  \item The gateway's connection pool to course-service fills; its
  \code{/api/courses} route degrades for everyone, including requests
  that never needed auth-service. Chapter~31's cascade begins one hop
  up.
  \item auth-service, meanwhile, is receiving 100 requests per second
  from a caller that will discard the answers after the gateway's
  timeout --- load that makes the vacuum's contention worse.
  \item The breaker never opens: 2\,s is under the 60\,s slow-call
  default, and none of the calls fail. The one mechanism that could
  shed load is watching for the wrong signal.
\end{enumerate}

With the configuration below: calls are cut at 2\,s (Feign read
timeout) and counted as slow above 1\,s; after ten outcomes with
$\ge$80\,\% slow, the breaker opens; listings return with the cached or
degraded teacher name in under a millisecond; at most 25 threads were
ever waiting; auth-service's load drops to three probe calls per ten
seconds; and the breaker's state change is on a dashboard.

\section{Enhancing the resolution}
\label{sec:res-enhance}

Dependency first: add \code{spring-cloud-starter-circuitbreaker-resilience4j}
to course-service's POM (it brings the \code{CircuitBreakerFactory}
Feign needs; \code{resilience4j-spring-boot3} stays for the annotations
and metrics). Then \code{course-service.yml}:

\begin{lstlisting}[language=yaml]
spring:
  cloud:
    openfeign:
      circuitbreaker:
        enabled: true                            (*@\co{1}@*)
      client:
        config:
          auth-service:                          (*@\co{2}@*)
            connect-timeout: 1000
            read-timeout: 2000
            logger-level: basic

resilience4j:
  circuitbreaker:
    instances:
      authService:
        sliding-window-type: COUNT_BASED
        sliding-window-size: 10
        minimum-number-of-calls: 10
        failure-rate-threshold: 50
        slow-call-duration-threshold: 1s         (*@\co{3}@*)
        slow-call-rate-threshold: 80
        wait-duration-in-open-state: 10s
        permitted-number-of-calls-in-half-open-state: 3
        automatic-transition-from-open-to-half-open-enabled: true
        register-health-indicator: true          (*@\co{4}@*)
        ignore-exceptions:                       (*@\co{5}@*)
          - feign.FeignException.NotFound
          - feign.FeignException.BadRequest
  retry:
    instances:
      authService:
        max-attempts: 3                          (*@\co{6}@*)
        wait-duration: 200ms
        enable-exponential-backoff: true
        exponential-backoff-multiplier: 2
        enable-randomized-wait: true
        randomized-wait-factor: 0.5
        retry-exceptions:
          - java.net.SocketTimeoutException
          - feign.RetryableException
        ignore-exceptions:
          - feign.FeignException.FeignClientException
  bulkhead:
    instances:
      authService:
        max-concurrent-calls: 25                 (*@\co{7}@*)
        max-wait-duration: 0
  timelimiter:
    instances:
      authService:
        timeout-duration: 3s                     (*@\co{8}@*)
        cancel-running-future: true

management:
  endpoints:
    web:
      exposure:
        include: health,info,metrics,prometheus,circuitbreakers,circuitbreakerevents
  health:
    circuitbreakers:
      enabled: true
\end{lstlisting}

\begin{description}
  \item[\co{1}] The switch that makes \code{fallback = ...} real. Feign
  now wraps each method in a breaker named
  \code{AuthServiceClient\#getUserById(Long)}; the explicit annotations
  below use the \code{authService} instance for the caller-side policy.
  \item[\co{2}] Per-client timeouts, keyed by the Feign \code{name}.
  Two seconds is generous for a by-id lookup that normally takes
  20\,ms, and strictly shorter than the gateway's 5\,s.
  \item[\co{3}] Slow now counts. Above 1\,s is ``slower than designed'';
  when 80\,\% of the window is slow, the breaker opens even if every
  call eventually succeeds.
  \item[\co{4}] The breaker's state appears in \code{/actuator/health}
  and as a Micrometer gauge (\code{circuitbreaker.state}) --- the thing
  you alert on.
  \item[\co{5}] A missing user is an answer, not an outage. 4xx
  responses neither count as failures nor trigger the fallback.
  \item[\co{6}] Three attempts with 200\,ms, then 400\,ms, each
  jittered by $\pm$50\,\%. Only transport failures are retried; any
  HTTP 4xx is final. Total worst case: $3 \times 2\,\text{s} + 0.6\,\text{s}$,
  still bounded and still under a minute --- but longer than the
  gateway's 5\,s, which is the argument for at most 2 attempts on a
  request path and 3 only in background work.
  \item[\co{7}] Twenty-five threads may wait on auth-service; the 26th
  call fails immediately with \code{BulkheadFullException} and takes the
  fallback. \code{max-wait-duration: 0} means no queueing --- queueing
  is what we are trying to avoid.
  \item[\co{8}] Only used by the \code{CompletableFuture} variant
  below; harmless otherwise.
\end{description}

The service-side wrapper, with identity propagation and a fallback that
says what it is:

\begin{lstlisting}[language=Java]
/** Forwards the caller's identity on service-to-service calls so
 *  auth-service's GatewayHeaderAuthFilter accepts them. */
@Configuration
public class FeignIdentityConfig {
    @Bean
    public RequestInterceptor identityPropagation() {
        return template -> {
            var auth = SecurityContextHolder.getContext()
                    .getAuthentication();
            if (auth != null && auth.getPrincipal() instanceof Long id) {
                template.header(JwtConstants.HEADER_USER_ID, id.toString());
                auth.getAuthorities().stream().findFirst().ifPresent(a ->
                    template.header(JwtConstants.HEADER_USER_ROLE,
                        a.getAuthority().replace("ROLE_", "")));
            }
        };
    }
}

/** The one place course-service asks who a user is. */
@Service
@RequiredArgsConstructor
public class UserLookupService {

    private final AuthServiceClient authServiceClient;
    private final Map<Long, UserDto> lastKnown =
            new ConcurrentHashMap<>();                     (*@\co{1}@*)

    @Retry(name = "authService")
    @CircuitBreaker(name = "authService", fallbackMethod = "degraded")
    @Bulkhead(name = "authService")
    public UserDto findUser(Long id) {
        UserDto user = authServiceClient.getUserById(id);
        lastKnown.put(id, user);
        return user;
    }

    /** Same signature plus the cause; called for breaker rejections,
     *  bulkhead rejections, timeouts and 5xx -- never for 4xx. */
    private UserDto degraded(Long id, Throwable cause) {   (*@\co{2}@*)
        UserDto cached = lastKnown.get(id);
        if (cached != null) {
            return cached;                                 // stale, real
        }
        return new UserDto(id, null, "Unknown", "user", null);
    }
}
\end{lstlisting}

\begin{description}
  \item[\co{1}] A last-known-good cache. In one replica a map is
  enough; with several, the same idea belongs in Redis with a TTL, or
  --- the real fix --- in a local table fed by user events.
  \item[\co{2}] The fallback is typed and visible: callers get a
  \code{UserDto} whose name says ``Unknown'', never \code{null}. The
  \code{Throwable} lets you log \emph{why} (breaker open vs.\ timeout)
  at WARN once per state change rather than per call.
\end{description}

The metrics to watch, all exported to Prometheus by the existing
Actuator setup: \code{resilience4j\_circuitbreaker\_state} (a gauge per
state per instance; alert when \code{state="open"} is 1 for more than a
minute), \code{resilience4j\_circuitbreaker\_calls\_seconds} by
\code{kind} (successful, failed, ignored) for rates,
\code{resilience4j\_bulkhead\_available\_concurrent\_calls} approaching
zero as the leading indicator, and Feign's own
\code{http\_client\_requests\_seconds} for the raw latency that started
it all.

\begin{decision}[when the call should not exist]
Keep the synchronous call while it is off the hot path, rare, and
tolerates a stale answer. Replace it with a local copy fed by events
when a listing page needs it (latency and coupling), when grading must
work while auth-service is down (availability), or when the two teams
release on different cadences (contract drift). Replace it with a
short-TTL cache in front of the call when the data is hot but must not
be more than a minute old and eventual consistency via events is not
worth a new consumer. The senior instinct is to count synchronous
dependencies per request and treat each one as a multiplication of
failure probability: three 99.9\,\% dependencies in series are a
99.7\,\% endpoint.
\end{decision}

\begin{handson}[watch the breaker open]
With the enhancement applied and the stack up, stop auth-service and
call an endpoint that uses \code{UserLookupService} twelve times:
\begin{lstlisting}[language=bash]
$ docker compose stop auth-service
$ for i in $(seq 1 12); do
    curl -s -o /dev/null -w "%{http_code} %{time_total}s\n" \
      -H "Authorization: Bearer $TOKEN" \
      http://localhost:8080/api/courses/courses/1
  done
200 1.04s   ... (10 slow lines: retries then fallback)
200 0.01s   <- breaker open: fallback without a network call
200 0.01s
$ curl -s http://localhost:8082/actuator/circuitbreakers | jq
{ "circuitBreakers": { "authService": {
    "failureRate": "100.0%", "slowCallRate": "0.0%",
    "bufferedCalls": 10, "failedCalls": 10, "state": "OPEN" } } }
\end{lstlisting}
Every response is 200 with ``Unknown user'' as the teacher --- the page
degrades instead of failing. Start auth-service again, wait 10\,s, and
\code{/actuator/circuitbreakerevents} shows the three HALF\_OPEN probes
succeed and the state return to CLOSED.
\end{handson}

\section{Outside the project}

The same machinery under other names: Hystrix was Netflix's original
and is retired; Resilience4j is its successor in the Java world;
Envoy and Istio implement timeouts, retries with budgets, and outlier
ejection in the sidecar so that application code contains none of it,
which is the direction large fleets have gone. gRPC carries deadlines
in the protocol --- a caller's deadline propagates to every hop, solving
the ``timeouts stack'' problem by construction. Cloud SDKs (AWS, GCP)
ship with exponential backoff and jitter built in for the same reason
this chapter recommends them. The vocabulary is universal; an
interviewer who says ``bulkhead'' or ``retry budget'' means exactly
what this chapter means.

\section{Questions from real interviews}

\subsection*{Q: auth-service latency goes from 20\,ms to 2\,s. Walk me through what happens in course-service, and how you would stop it.}

Each request that touches auth-service now holds a Tomcat thread a
hundred times longer. At a modest rate the pool of 200 saturates within
seconds, the accept queue fills, and course-service stops answering
everything --- including endpoints that never call auth-service and the
readiness probe. Upstream, the gateway's per-host pool to course-service
fills and its route degrades for all users. Nothing in the default
configuration stops this: Feign waits 60\,s, and a 2\,s call is not
slow to a breaker whose slow-call threshold is 60\,s. To stop it, in
order of leverage: a read timeout of 2\,s so no thread waits longer
than the gateway's deadline; a bulkhead of about 25 so the other 175
threads keep serving; a slow-call threshold of 1\,s at 80\,\% so the
breaker opens on latency, not just on errors; and a fallback that
returns the last known name rather than \code{null}. Then, if the
lookup is on a hot path, remove the call altogether with a local copy
fed by user events. The senior framing: the goal is that auth-service's
latency shows up as a stale name on a page, not as an outage of
courses.

\subsection*{Q: How do you size timeouts across three hops?}

From the outside in, each shorter than the one before. The browser or
mobile client is the outermost budget --- say 10\,s before the user
gives up. The gateway's response timeout must be inside that, 5\,s,
so the client gets a clean 503 instead of a hang. course-service's
Feign read timeout must be inside \emph{that} and leave room for
retries: 2\,s with at most one retry keeps the worst case around 4\,s.
Any timeout longer than its caller's is a leak --- work done for a
request that has already been abandoned --- and any retry policy must
be included in the arithmetic, because three attempts at 2\,s is 6\,s
plus backoff. In production you also set a \emph{connect} timeout
separately and much shorter (1\,s), because failure to connect is
diagnosed in one round trip. The best systems propagate a deadline
header from the edge so each hop computes its own remaining budget;
gRPC does this natively.

\subsection*{Q: Retry storms --- how do you prevent them?}

Three controls, all needed. First, jitter: exponential backoff alone
lets every client that failed at the same moment retry at the same
moment, so the second wave is synchronized and larger; randomizing the
wait by $\pm$50\,\% spreads it. Second, a retry budget: retries as a
percentage of traffic (Envoy's default is 20\,\%), so that when a
dependency is truly down, retries cannot triple its load. Resilience4j
approximates this by putting Retry outside the breaker --- an open
breaker rejects attempts before they become retries. Third,
idempotency: retry only \code{GET} and operations that carry an
idempotency key the server deduplicates on; otherwise a timeout on a
\code{POST} becomes a duplicate enrollment. And retry only transport
failures and 503 --- a 400 or 404 is deterministic and will fail again.
The tell-tale sign of a storm in metrics is a dependency whose request
rate rises as its success rate falls.

\subsection*{Q: The breaker is open, but users need the data. What does your fallback return?}

The most truthful thing available, typed so the caller cannot mistake
it for the real value. Order of preference: the last successful value
from a short-lived cache (stale but real --- a teacher's name from ten
minutes ago is fine); a clearly degraded placeholder the UI can render
(``Unknown teacher'') that never masquerades as data; an empty
collection where the domain permits; or a typed exception that becomes
a 503 with \code{Retry-After}. Never \code{null}, which turns a
resilience event into a \code{NullPointerException} somewhere else. The
fallback must also see the cause: a 404 means the user does not exist
and should surface as such, not be swallowed by ``try later''. And
fallbacks should be cheap and local --- a fallback that calls another
service is a second dependency on the failure path.

\subsection*{Q: How do you know the breaker is open at 3\,a.m.?}

Because it is a metric with an alert, not a log line. Resilience4j
exports \code{resilience4j\_circuitbreaker\_state} as a gauge per
state; a rule like ``\code{state="open"} equals 1 for more than 60
seconds'' pages someone with the dependency name in the alert. Pair it
with the leading indicators: \code{slow\_call\_rate} rising before the
breaker trips, bulkhead available slots trending to zero, and the
dependency's p99 from \code{http\_client\_requests\_seconds}. The
breaker also feeds \code{/actuator/health} when
\code{register-health-indicator} is on, so a readiness probe can be
made to reflect it --- though usually you do \emph{not} want an open
breaker to take a service out of rotation, since it is serving
fallbacks precisely so it can stay in. And log the \emph{transition}
at WARN once, not every rejected call, or the log becomes the next
outage.

\subsection*{Q: Would you put the breaker in the gateway or in the service?}

Both, for different failures. The gateway's breaker (Chapter~31)
protects users from a whole service being slow; it opens per route and
serves a fallback body. The service's breaker protects the service's
own thread pool from one of its dependencies and is where the
domain-aware fallback lives --- only course-service knows that a
missing teacher name is survivable. A gateway cannot know that. The
mistake is having only the edge one and assuming services behind it
are therefore safe; the cascade this chapter walks through happens
entirely behind the gateway.

\begin{quiz}
\begin{enumerate}
  \item The project configures a breaker named \code{authService}.
  List the three separate reasons it currently protects no call, and
  the one-line change that fixes each.
  \item With \code{sliding-window-size: 10} and the other defaults,
  how many consecutive failures can occur after startup before the
  breaker opens, and why does a call that takes 30\,s not count?
  \item Order these and justify: gateway response timeout, browser
  timeout, Feign read timeout, retry count.
  \item \code{AuthServiceFallback.getUserById} returns \code{null}. Name
  two concrete defects that follow from this and the two-argument
  signature that fixes the more subtle one.
  \item auth-service returns 401 to every Feign call. Explain the cause
  in the project's security configuration, the effect on the breaker,
  and the two changes that resolve it.
  \item Three synchronous dependencies each at 99.9\,\% availability
  sit in series on one endpoint. What is its availability, and which
  two design changes raise it without touching the dependencies?
\end{enumerate}
\end{quiz}
